Segmenting developing bones in pediatric hand radiographs is difficult when adjacent ossification centers are separated by narrow, low-contrast gaps. Region-overlap objectives can produce high Dice scores while still joining structures that should remain separate. We introduce instance-derived seam supervision, a training-only construction that converts indexed bone annotations into a narrow background class between neighboring instances. The resulting three-class target can be used with a conventional semantic segmentation backbone and requires only the radiograph at inference. We evaluate the construction with matched nnU-Net models on 875 training and 96 validation radiographs under a frozen evaluation protocol. Relative to a mature binary nnU-Net, seam supervision reduces the image-level component under-count rate from 51.04\% to 42.71\% and the gap-region false-positive rate from 0.1873 to 0.1832. Dice remains nearly unchanged (0.8980 versus 0.8974), while Boundary IoU and 2-pixel Surface Dice increase by 0.0007 and 0.0015. The gains are therefore specific but modest: HD95 and ASSD worsen, revealing occasional boundary outliers. These results show that instance-derived background supervision can reduce adjacent-structure merging, but also motivate confidence weighting and foreground-preservation constraints before the approach is claimed as a general plug-in module.
